% =====================================================================
%  Предзащита ВКР — Прогнозирование ВРП регионов РФ
%  XeLaTeX + Beamer 16:9 | HSE style
% =====================================================================
\documentclass[aspectratio=169,10pt]{beamer}

% ---------- Пакеты -------------------------------------------------
\usepackage{fontspec}
\usepackage{polyglossia}
\setmainlanguage{russian}
\setotherlanguage{english}
\setmainfont{DejaVu Serif}
\setsansfont{DejaVu Sans}
\setmonofont{DejaVu Sans Mono}

\usepackage{graphicx}
\usepackage{booktabs}
\usepackage{xcolor}
\usepackage{tikz}
\usepackage{amsmath}

% ---------- Цвета HSE ----------------------------------------------
\definecolor{hseblue}{HTML}{0B3D91}
\definecolor{hselightblue}{HTML}{5B8DB8}
\definecolor{hsegray}{HTML}{F2F2F2}
\definecolor{hsedarkgray}{HTML}{4A4A4A}
\definecolor{hseaccent}{HTML}{D4312C}

% ---------- Beamer тема --------------------------------------------
\usetheme{default}
\usecolortheme{default}
\usefonttheme{professionalfonts}

\setbeamercolor{frametitle}{fg=white, bg=hseblue}
\setbeamercolor{title}{fg=white}
\setbeamercolor{subtitle}{fg=hselightblue!80!white}
\setbeamercolor{author}{fg=white}
\setbeamercolor{date}{fg=hselightblue!70!white}
\setbeamercolor{institute}{fg=hselightblue!70!white}
\setbeamercolor{titlelike}{fg=hseblue}
\setbeamercolor{structure}{fg=hseblue}
\setbeamercolor{normal text}{fg=hsedarkgray}
\setbeamercolor{block title}{bg=hseblue, fg=white}
\setbeamercolor{block body}{bg=hsegray, fg=hsedarkgray}
\setbeamercolor{alerted text}{fg=hseaccent}

\setbeamertemplate{navigation symbols}{}
\setbeamertemplate{frametitle}{
  \begin{beamercolorbox}[wd=\paperwidth, dp=0.18cm, leftskip=0.6cm, rightskip=0.4cm, sep=0.16cm]{frametitle}
    \usebeamerfont{frametitle}\insertframetitle
  \end{beamercolorbox}
}

\setbeamertemplate{footline}{
  \begin{beamercolorbox}[wd=\paperwidth, ht=0.4cm, dp=0.15cm]{structure}
    \hfill
    \textcolor{hseblue!60}{\scriptsize \insertframenumber\,/\,\inserttotalframenumber}
    \hspace{0.4cm}
  \end{beamercolorbox}
}

\setbeamertemplate{itemize item}{\textcolor{hseblue}{\small$\blacktriangleright$}}
\setbeamertemplate{itemize subitem}{\textcolor{hselightblue}{\small$\bullet$}}

% ---------- Титульный слайд ----------------------------------------
\title{\large Машинное обучение в прогнозировании\\валового регионального продукта\\субъектов Российской Федерации}
\subtitle{Предзащита выпускной квалификационной работы}
\author{Семенюченко Артём Геннадьевич}
\institute{Научный руководитель: Паточенко Евгений Анатольевич\\
           Программа «Анализ больших данных», ФКН НИУ ВШЭ}
\date{2026}

% ===================================================================
\begin{document}

% ------ Слайд 1: Титульный ----------------------------------------
\begin{frame}[plain]
\begin{tikzpicture}[remember picture, overlay]
  \fill[hseblue] (current page.north west) rectangle (current page.south east);
  \fill[hseblue!80!black] (current page.south west) rectangle
        ([yshift=0.9cm]current page.south east);
\end{tikzpicture}
\vspace{0.9cm}
\titlepage
\end{frame}


% ------ Слайд 2: Мотивация ----------------------------------------
\begin{frame}{Мотивация и актуальность}
\begin{columns}[T]
\column{0.55\textwidth}
\textbf{Почему это важно?}
\medskip
\begin{itemize}\setlength\itemsep{4pt}\small
  \item ВРП — ключевой показатель социально-экономического развития регионов
  \item Используется для межбюджетных трансфертов и регионального планирования
  \item Официальные данные Росстат публикуются с лагом 1–2 года
  \item Традиционные методы (ARIMA, VAR) не учитывают нелинейности и пространственную гетерогенность
\end{itemize}
\column{0.42\textwidth}
\begin{block}{Исследовательский вопрос}
\small Могут ли методы ML обеспечить точный прогноз ВРП для 85 регионов РФ на горизонте 1 год, превосходя наивные бенчмарки?
\end{block}
\smallskip
\begin{block}{Практическая ценность}
\small Ранняя оценка ВРП позволяет оперативно корректировать политику регионального развития
\end{block}
\end{columns}
\end{frame}


% ------ Слайд 3: Цель и задачи ------------------------------------
\begin{frame}{Цель и задачи исследования}
\begin{columns}[T]
\column{0.48\textwidth}
\begin{alertblock}{Цель}
\small Разработать ML-пайплайн прогнозирования ВРП регионов РФ с интерпретируемостью через SHAP
\end{alertblock}
\medskip
\textbf{Задачи:}
\begin{enumerate}\setlength\itemsep{3pt}\small
  \item Панельный датасет 2005–2023, 85 регионов
  \item Расширяющееся окно без утечки данных
  \item Сравнение бустинга (LGBM, CatBoost, GPBoost) с бенчмарками
  \item Аблационный анализ блоков признаков
  \item Устойчивость к выбросам и кризисам
  \item Глобальный и локальный SHAP-анализ
\end{enumerate}
\column{0.48\textwidth}
\begin{block}{Новизна работы}
\begin{itemize}\setlength\itemsep{4pt}
  \item Первое систематическое сравнение ML-моделей на панельных данных ВРП всех регионов РФ
  \item Интеграция GPBoost (гауссовы процессы + бустинг) для учёта пространственной корреляции
  \item Блочный SHAP-анализ по типам регионов
\end{itemize}
\end{block}
\end{columns}
\end{frame}


% ------ Слайд 4: Данные -------------------------------------------
\begin{frame}{Данные}
\begin{columns}[T]
\column{0.55\textwidth}
\textbf{Источник:} Росстат, база данных показателей субъектов РФ
\medskip
\begin{itemize}\setlength\itemsep{4pt}
  \item \textbf{Период:} 2005–2023 гг.
  \item \textbf{Объектов:} 85 регионов (отфильтровано $\geq$80 по покрытию)
  \item \textbf{Целевая переменная:} $\Delta \log(\text{ВРП}_{\text{реал}})$ — темп прироста реального ВРП
  \item \textbf{Блоки признаков:}
  \begin{itemize}
    \item Лаги ВРП и темпов роста (AR-компонента)
    \item Инвестиции, доходы, промышленность
    \item Занятость, демография
    \item Бюджетные показатели
    \item Временны́е и региональные эффекты
  \end{itemize}
\end{itemize}
\column{0.42\textwidth}
\begin{block}{Предобработка}
\begin{itemize}\setlength\itemsep{3pt}
  \item Логарифмирование стоимостных показателей
  \item RobustScaler по train-выборке каждого фолда
  \item Медианная импутация пропусков
  \item Порог пропусков: 40\%
  \item Исключение признаков с будущей информацией
\end{itemize}
\end{block}
\begin{block}{Кризисные периоды}
\small 2008–09, 2014–15, 2020, 2022–23
\end{block}
\end{columns}
\end{frame}


% ------ Слайд 5: Методология --------------------------------------
\begin{frame}{Методология: расширяющееся скользящее окно}
\begin{columns}[T]
\column{0.54\textwidth}
\textbf{Схема валидации (9 фолдов):}
\medskip

\begin{tabular}{lcc}
\toprule
Фолд & Train & Test \\
\midrule
1 & 2005–2014 & 2015 \\
2 & 2005–2015 & 2016 \\
$\vdots$ & $\vdots$ & $\vdots$ \\
9 & 2005–2022 & 2023 \\
\bottomrule
\end{tabular}
\medskip

\begin{itemize}\setlength\itemsep{4pt}
  \item Минимум 10 лет на обучение
  \item \textbf{FeatureBuilder} обучается только на train
  \item Нет утечки через масштабирование и импутацию
\end{itemize}

\column{0.43\textwidth}
\begin{block}{Метрики}
$$\text{RMSE} = \sqrt{\frac{1}{N}\sum_i(y_i - \hat{y}_i)^2}$$
$$\text{MAE} = \frac{1}{N}\sum_i|y_i - \hat{y}_i|$$
$$\text{OOS-}R^2 = 1 - \frac{\sum(y_i-\hat{y}_i)^2}{\sum(y_i-\bar{y}_{\text{train}})^2}$$
\end{block}
\medskip
\begin{alertblock}{Бенчмарк}
Наивная модель: $\hat{y}_i = \bar{y}_{\text{train}}$ (среднее темпа роста)
\end{alertblock}
\end{columns}
\end{frame}


% ------ Слайд 6: Инженерия признаков -----------------------------
\begin{frame}{Архитектура признаков}
\begin{columns}[T]
\column{0.48\textwidth}
\textbf{Блоки признаков (аблационный анализ):}
\medskip
\begin{itemize}\setlength\itemsep{5pt}
  \item \textbf{A0 — Полный набор} (baseline)
  \item \textbf{A1 — без AR-лагов}: влияние автокорреляции ВРП
  \item \textbf{A2 — без инвестиций}: роль капитальных вложений
  \item \textbf{A3 — без доходов}: эффект доходов населения
  \item \textbf{A4 — без бюджета}: бюджетный канал
  \item \textbf{A5 — без временны́х эффектов}: макросинхронизация
\end{itemize}
\column{0.48\textwidth}
\begin{block}{Принцип «no leakage»}
Все преобразования признаков (скейлинг, импутация, кодирование) применяются строго на train-части каждого фолда. Test-данные преобразуются параметрами, оцененными только на train.
\end{block}
\medskip
\begin{block}{Лаговая структура}
Для цели $\Delta\log(\text{ВРП}_t)$ используются значения признаков $t-1$, что исключает информацию о тестовом периоде.
\end{block}
\end{columns}
\end{frame}


% ------ Слайд 7: Модели -------------------------------------------
\begin{frame}{Модели прогнозирования}
\begin{columns}[T]
\column{0.48\textwidth}
\textbf{ML-модели:}
\begin{itemize}\setlength\itemsep{4pt}
  \item \textbf{LGBM} — LightGBM, gradient boosting на деревьях
  \item \textbf{CatBoost} — Yandex, нативная обработка категориальных признаков
  \item \textbf{GPBoost} — комбинация GB + гауссовые процессы, учёт случайных эффектов регионов
\end{itemize}
\medskip
\textbf{Бенчмарки:}
\begin{itemize}\setlength\itemsep{4pt}
  \item \textbf{Ridge} — L2-регрессия, перекрёстная валидация $\lambda$
  \item \textbf{AR(1)} — авторегрессия в пространстве темпов роста
  \item \textbf{Naive} — среднее темпа роста по train
\end{itemize}
\column{0.48\textwidth}
\begin{block}{Подбор гиперпараметров}
\begin{itemize}\setlength\itemsep{2pt}\small
  \item Optuna TPE, 20 итераций
  \item Быстрые trial: n\_estimators снижен
  \item Финальный fit: полные 300/500/300
  \item Inner-val: последние 2 года train
\end{itemize}
\end{block}
\smallskip
\begin{block}{Интерпретируемость}
\small \texttt{shap.TreeExplainer}: глобальный, групповой и локальный SHAP
\end{block}
\end{columns}
\end{frame}


% ------ Слайд 8: Основные результаты — таблица -------------------
\begin{frame}{Основные результаты: сводная таблица метрик}
\vspace{-6pt}
\begin{center}
\includegraphics[width=0.92\textwidth,trim=0 30 0 30,clip]{grp_forecast/figures/metrics_table.png}
\end{center}
\vspace{-2pt}
\centering\small Зелёным — лучшее значение. OOS-$R^2>0$: модель превосходит наивный бенчмарк.
\end{frame}


% ------ Слайд 9: Динамика RMSE по фолдам -------------------------
\begin{frame}{Динамика RMSE по годам прогноза}
\begin{center}
\includegraphics[width=0.90\textwidth]{grp_forecast/figures/rmse_over_time.png}
\end{center}
\vspace{-4pt}
\small
Серые полосы — кризисные периоды (2008–09, 2014–15, 2020, 2022–23).
Все модели показывают всплески ошибок в кризисные годы, что ожидаемо.
\end{frame}


% ------ Слайд 10: Ошибки по регионам -----------------------------
\begin{frame}{Ошибки прогноза: региональная структура}
\begin{columns}[T]
\column{0.50\textwidth}
\centering
\textbf{Топ-20 регионов по MAE}\\[4pt]
\includegraphics[width=\textwidth]{grp_forecast/figures/error_by_region.png}
\column{0.48\textwidth}
\centering
\textbf{Тепловая карта ошибок (регион × год)}\\[4pt]
\includegraphics[width=\textwidth]{grp_forecast/figures/error_heatmap.png}
\end{columns}
\vspace{2pt}
\small Наибольшие ошибки — в ресурсодобывающих регионах (ЯНАО, Сахалин) и в периоды структурных шоков.
\end{frame}


% ------ Слайд 11: SHAP глобальный ---------------------------------
\begin{frame}{SHAP: глобальная важность признаков}
\begin{columns}[T]
\column{0.49\textwidth}
\centering
\textbf{Средний |SHAP|}\\[4pt]
\includegraphics[width=\textwidth,height=5.2cm,keepaspectratio]{grp_forecast/figures/shap_global.png}
\column{0.49\textwidth}
\centering
\textbf{Beeswarm (направление эффекта)}\\[4pt]
\includegraphics[width=\textwidth,height=5.2cm,keepaspectratio]{grp_forecast/figures/shap_beeswarm.png}
\end{columns}
\vspace{2pt}
\small Доминирует AR-компонента; инвестиции и доходы — второй блок. Beeswarm подтверждает ожидаемые знаки эффектов.
\end{frame}


% ------ Слайд 12: SHAP групповой ----------------------------------
\begin{frame}{SHAP: блоки признаков по типам регионов}
\begin{center}
\includegraphics[width=0.80\textwidth]{grp_forecast/figures/shap_group_heatmap.png}
\end{center}
\vspace{-4pt}
\small
Тепловая карта показывает, что для ресурсодобывающих регионов важнее отраслевые и инвестиционные признаки,
а для регионов-лидеров сервисной экономики (Москва, СПб) — потребительские и бюджетные блоки.
\end{frame}


% ------ Слайд 13: Аблационный анализ ------------------------------
\begin{frame}{Аблационный анализ блоков признаков}
\begin{center}
\includegraphics[width=0.74\textwidth]{grp_forecast/figures/ablation_study.png}
\end{center}
\vspace{-2pt}
\small A0 — полный набор. Удаление AR-лагов (A1) и инвестиций (A2) даёт наибольший рост RMSE.
\end{frame}


% ------ Слайд 14: Устойчивость ------------------------------------
\begin{frame}{Устойчивость модели}
\begin{columns}[T]
\column{0.49\textwidth}
\centering
\textbf{Исключение структурных выбросов}\\[4pt]
\includegraphics[width=\textwidth]{grp_forecast/figures/rc1_outliers.png}
\small\\
Полная выборка vs. без Москвы, ХМАО, ЯНАО, Сахалина
\column{0.49\textwidth}
\centering
\textbf{Кризисные vs. обычные периоды}\\[4pt]
\includegraphics[width=\textwidth]{grp_forecast/figures/rc2_crisis.png}
\small\\
RMSE на фолдах кризисных лет существенно выше
\end{columns}
\vspace{4pt}
\small ML-модели устойчивы к исключению структурных выбросов: RMSE меняется несущественно.
На кризисных фолдах точность падает, однако ML превосходит бенчмарки и в этом режиме.
\end{frame}


% ------ Слайд 15: Заключение --------------------------------------
\begin{frame}{Заключение и вклад}
\begin{columns}[T]
\column{0.55\textwidth}
\textbf{Основные результаты:}
\begin{itemize}\setlength\itemsep{4pt}\small
  \item Лучшая модель (LGBM / CatBoost) превосходит наивный бенчмарк на всех 9 фолдах
  \item OOS-$R^2 > 0$ — информативное прогнозирование подтверждено
  \item AR-лаги и инвестиции — ключевые предикторы (SHAP + аблация)
  \item Устойчивость к выбросам; ошибки растут в кризисные периоды
\end{itemize}
\column{0.42\textwidth}
\begin{block}{Научный вклад}
\begin{itemize}\setlength\itemsep{3pt}\small
  \item Воспроизводимый пайплайн для панельного прогнозирования ВРП
  \item Строгая OOS-методология для региональной экономики
  \item Блочный SHAP — экономическая интерпретация ML
\end{itemize}
\end{block}
\smallskip
\begin{alertblock}{Практика}
\small Оперативная оценка ВРП для регионального планирования и МБТ
\end{alertblock}
\end{columns}
\medskip
\centering
\textcolor{hseblue}{\large Спасибо за внимание!}
\end{frame}


\end{document}
